% =========================================================
% Proof: Differentiability and One-Sided Derivatives
% Source: volume-iii/analysis/differentiation/notes/notes-onesided.tex
% Mode: standard
% =========================================================

\subsection*{Differentiability and One-Sided Derivatives}
\label{prf:differentiability-one-sided}

\begin{remark*}[Return]
\hyperref[thm:differentiability-one-sided]{$\leftarrow$ Back to Theorem (Differentiability and One-Sided Derivatives) in Notes}
\end{remark*}

\begin{proof}
Let $I \subseteq \mathbb{R}$ be an open interval and let $f : I \to \mathbb{R}$.
Fix $c \in I$. The claim is that $f$ is differentiable at $c$ if and only if
both one-sided derivatives $f'_-(c)$ and $f'_+(c)$ exist and are equal, in
which case $f'(c) = f'_-(c) = f'_+(c)$.

$(\Rightarrow)$ \textit{Assume $f$ is differentiable at $c$; deduce that
$f'_-(c)$ and $f'_+(c)$ exist and are equal.}

By \ByDef{Differentiability at a Point}, the two-sided limit
\[
  \lim_{x \to c} \frac{f(x) - f(c)}{x - c} = f'(c)
\]
exists. By \ByThm{Two-Sided Limit and One-Sided Limits}, a two-sided limit
exists if and only if both one-sided limits exist and agree. Therefore
$f'_-(c)$ and $f'_+(c)$ both exist and satisfy
$f'_-(c) = f'_+(c) = f'(c)$.

$(\Leftarrow)$ \textit{Assume $f'_-(c)$ and $f'_+(c)$ both exist and are
equal to some $L \in \mathbb{R}$; deduce that $f$ is differentiable at $c$.}

By hypothesis,
\[
  \lim_{x \to c^-} \frac{f(x) - f(c)}{x - c} = L
  \qquad \text{and} \qquad
  \lim_{x \to c^+} \frac{f(x) - f(c)}{x - c} = L.
\]
Since both one-sided limits of the difference quotient exist and are equal,
the two-sided limit exists and equals $L$ by \ByThm{Two-Sided Limit and
One-Sided Limits}. By \ByDef{Differentiability at a Point}, $f$ is
differentiable at $c$ and $f'(c) = L$.
\end{proof}

\begin{remark*}[Proof shape]
Both directions reduce immediately to the characterisation of a two-sided
limit in terms of its one-sided limits: the derivative is defined as a
two-sided limit of the difference quotient, so the equivalence is inherited
directly from that limit-theoretic fact. The proof introduces no construction
and requires no calculation — it is a pure restatement via the one-sided
limit criterion.

\FlashProof{1. ($\Rightarrow$) The derivative is a two-sided limit; by the
  one-sided limit criterion, both $f'_-(c)$ and $f'_+(c)$ exist and equal
  $f'(c)$.
  2. ($\Leftarrow$) Both one-sided limits of the difference quotient equal
  $L$; by the one-sided limit criterion the two-sided limit exists, so $f$
  is differentiable with $f'(c) = L$.}
\end{remark*}

\begin{remark*}[Dependencies]
\begin{itemize}
  \item \hyperref[def:derivative-at-point]{Definition of Differentiability
    at a Point}: Identifies the derivative as the two-sided limit of the
    difference quotient, in both directions.
  \item \hyperref[thm:two-sided-one-sided-limits]{Two-Sided Limit and
    One-Sided Limits}: The single result doing all the work — a two-sided
    limit exists iff both one-sided limits exist and agree.
\end{itemize}

\FlashUses{def:derivative-at-point, thm:two-sided-one-sided-limits}
\end{remark*}